\documentclass[12pt]{article}
\usepackage[T1]{fontenc}
\usepackage[utf8]{inputenc}
\usepackage{lmodern}
\usepackage{textcomp}
\usepackage{enumitem}
\usepackage{geometry}
\geometry{margin=1in}
\begin{document}
\begin{center}
Answer Key - Biology - Graduate Level Assessment 19 \
\small (Answers are ordered exactly as in the question paper.)
\end{center}

\section*{Multiple Choice}
\begin{enumerate}[label=\arabic*.]
\item \textbf{Answer:} D
\\ \textit{Options:} A) Guttation is the outcome of rapid cellular respiration in plant tissues.; B) Guttation is a result of photosynthesis.; C) Guttation is caused by external weather conditions.; D) Guttation is explained by the effect of root pressure.; E) Guttation is a byproduct of the plant's defense mechanisms.
\\ \textit{Note:} Guttation occurs when excess water in the plant roots creates pressure, forcing water out through specialized glands in the leaves, a phenomenon primarily driven by root pressure.
\item \textbf{Answer:} B
\\ \textit{Options:} A) The absence of a Y chromosome determines maleness in mammals; B) The presence of a Y chromosome usually determines maleness in mammals; C) The ratio of two X chromosomes to one Y chromosome determines maleness; D) The presence of an X chromosome determines maleness in mammals; E) The presence of an X chromosome in the absence of a Y chromosome determines maleness
\\ \textit{Note:} The correct answer is supported by the key fact that the presence of the Y chromosome carries the SRY gene, which initiates male sex determination in mammals, establishing a clear link between the Y chromosome and maleness.
\item \textbf{Answer:} A
\\ \textit{Options:} A) The stomach has a high permeability to hydrogen ions; B) The presence of beneficial bacteria that neutralize stomach acid; C) The stomach lining is replaced every month; D) The stomach lining secretes digestive enzymes only when food is present; E) The production of bile in the stomach neutralizes the digestive acids
\\ \textit{Note:} The correct answer is based on the fact that the stomach lining produces a thick layer of mucus and bicarbonate that protects it from the corrosive effects of gastric acid and prevents self-digestion, rather than its permeability to hydrogen ions.
\item \textbf{Answer:} A
\\ \textit{Options:} A) The zygote will not divide into two embryos, producing twins.; B) The sperm will not be able to fertilize the egg.; C) The sperm will fertilize the ovum.; D) The sperm will fertilize multiple eggs.; E) The zygote will begin to divide normally.
\\ \textit{Note:} The acrosome and cortical reactions prevent polyspermy, ensuring that only one sperm fertilizes the egg, which is critical for the proper development of a single zygote and avoids the formation of multiple embryos.
\item \textbf{Answer:} A
\\ \textit{Options:} A) It attracts predators away from the nestlings.; B) It aids in temperature regulation of the nest by reducing insulation.; C) It assists in the hatching of other eggs.; D) It is a ritualistic behavior without any particular significance.; E) It signals to other gulls that the nesting site is occupied.
\\ \textit{Note:} The behavior of black-headed gulls removing broken eggshells from their nests serves the biological significance of reducing the likelihood of attracting predators to the nestlings, as the discarded shells can signal the presence of vulnerable young birds to potential threats.
\end{enumerate}

\section*{True / False}
\begin{enumerate}[label=\arabic*.]
\setcounter{enumi}{5}
\item \textbf{Answer:} False
\\ \textit{Options:} A) True; B) False
\\ \textit{Note:} The metastasis rate of station 9 lymph nodes was significantly lower than other mediastinal stations in lung cancer patients. The metastasis status of station 9 had no significant influence on tumor staging or prognosis. Routine dissection of station 9 lymph nodes may not be necessary, especially in patients with a low T stage, upper or middle lobe tumors, or without intrapulmonary lymph node metastasis.
\item \textbf{Answer:} True
\\ \textit{Options:} A) True; B) False
\\ \textit{Note:} Our data, derived from patients with coronary artery disease, support the hypothesis regarding a possible preventive effect of bezafibrate on the development of colon cancer.
\item \textbf{Answer:} False
\\ \textit{Options:} A) True; B) False
\\ \textit{Note:} There is a knowledge deficit regarding the SU treatment goal among gout patients receiving ULT, despite generally high levels of other gout-specific knowledge. SU goal information may be an important and underutilized concept among providers treating gout patients.
\item \textbf{Answer:} False
\\ \textit{Options:} A) True; B) False
\\ \textit{Note:} Results of this survey highlight the wide variety of practice patterns in the US for handling surgical margins in breast-conservation treatment. This issue remains controversial, with no prevailing standard of care. Consequently, additional study is needed in the modern era of multimodality treatment to examine the minimal amount of surgical treatment necessary, in conjunction with chemotherapy and radiation, to attain adequate local control rates in breast-conservation treatment.
\item \textbf{Answer:} False
\\ \textit{Options:} A) True; B) False
\\ \textit{Note:} This study did not demonstrate preventive effects of family meetings on the mental health of family caregivers. Further research should determine whether this intervention might be more beneficial if provided in a more concentrated dose, when applied for therapeutic purposes or targeted towards subgroups of caregivers.
\end{enumerate}

\section*{Long Form Response}
\begin{enumerate}[label=\arabic*.]
\setcounter{enumi}{10}
\item \textbf{Answer:} A newborn can acquire syphilis from an infected mother through two primary mechanisms: prenatal transmission and the birthing process. During the first four months of pregnancy, the bacterium Treponema pallidum can cross the placenta, leading to congenital syphilis. This transplacental transmission is particularly concerning as it can occur even in the absence of visible symptoms in the mother. Additionally, during the birthing process, if the fetus passes through the mother's infected vaginal canal, it can come into direct contact with syphilitic lesions, facilitating infection. Both routes underscore the critical need for prenatal screening and treatment of syphilis to prevent vertical transmission and protect neonatal health.
\\ \textit{Note:} correct because it accurately identifies and explains the two main routes of syphilis transmission from an infected mother to her newborn-transplacental transmission during pregnancy and exposure during the birthing process-emphasizing that the bacterium can cross the placenta even when the mother shows no symptoms.
\item \textbf{Answer:} Reduced salt intake plays a significant role in the regulation of blood pressure, particularly in individuals with hypertension. When salt (sodium) consumption decreases, the body retains less water due to the osmotic balance maintained by sodium levels, leading to a reduction in blood volume. This decrease in blood volume subsequently results in lower cardiac output and reduced pressure against the arterial walls, ultimately lowering blood pressure. Furthermore, the kidneys respond to lower sodium levels by increasing the excretion of sodium and water, enhancing this hypotensive effect. Overall, the physiological mechanisms of decreased salt intake effectively contribute to blood pressure regulation, providing a non-pharmacological approach to managing hypertension.
\\ \textit{Note:} correct because it accurately describes how reduced sodium intake decreases blood volume and cardiac output, leading to lower blood pressure through physiological mechanisms involving water retention and renal response.
\end{enumerate}

\vfill
\noindent\textit{End of Answer Key}
\end{document}